\BourbakiUnnumberedChapter{历史注记}{历史注记（第六、七章）}

（第六章和第七章）

（注：括号中的罗马数字指本文末尾的参考文献。）

初等算术运算，特别是分数的运算，不可避免地会引出许多关于整数之间整除性的经验观察。但无论是巴比伦人（尽管他们是代数学的专家），还是埃及人（尽管他们处理分数时技巧娴熟），似乎都未曾掌握支配这些性质的一般规则，而在这方面开创先河的是希腊人。他们的算术著作，其精妙阐述可见于欧几里得（I）的第七卷和第九卷，丝毫不逊色于他们在数学其他分支中最杰出的发现。两个整数的最大公约数的存在性，在第七卷的开头便通过被称为“欧几里得算法”* 的程序得以证明；它构成了后续所有发展的基础（素数的性质、最小公倍数的存在性与计算等），其论证方式与上文第六章第一节中的并无实质差异；而最辉煌的成就则是由两个卓越的定理构成，它们证明了素数的无穷性（第九卷命题20），并给出了构造偶完全数的方法（参见第七章第53页习题24；事实上，这一方法给出了所有偶完全数，正如欧拉后来所证明的那样）。唯有素数分解的存在性和唯一性未被以一般方式证明；然而，欧几里得确实明确证明了每个整数都能被某个素数整除（第七卷命题31），以及以下两个命题（第九卷命题13和14）：

“如果有任意多个数成连比例，首项为 1，公比恒定 {[}即等比数列{]}，且 1 后面的那个数是素数，那么最大的数将不能被任何数整除，除非这些数出现在该数列中”（换句话说，一个幂 \(p^n\) ——素数 \(p\) 的幂——只能被 \(p\) 的幂整除，而这些幂的指数 \(\leqslant n\)）。

“如果一个数是能被 {[}给定的{]} 素数整除的最小数，则它不能被任何其他素数整除，最初 {[}给定的{]} 整除它的素数除外”（换句话说，不同素数 \(p_{1}, \ldots, p_{k}\) 的乘积没有 \(p_{1}, \ldots, p_{k}\) 之外的素因子）。

因此看来，如果欧几里得没有陈述一般定理，那仅仅是因为缺乏适当的术语和记号来表示整数的任意次幂。*

尽管仔细研究使人觉得欧几里得的文本很可能由几个连续的层次组成，每一层对应于算术发展的一个阶段**，但似乎这一演变完全发生在公元前5世纪初至公元前4世纪中叶之间，人们不得不钦佩它所体现的技巧和逻辑自信：算术领域下一次可与之相比的进步要等到两千年之后才出现。

正是所谓的「不定」或「丢番图」问题构成了数论后续发展的根源。今天使用的「丢番图方程」一词在历史上并不完全恰当；它通常被理解为具有整数系数的多项式方程（或方程组），只寻求整数解：当方程是「确定的」，即只有有限多个（实数或复数）解时，这样的问题通常是不可能的；但另一方面，当未知数多于方程数时，它往往有解。现在，虽然丢番图确实似乎是第一个考虑「不定」问题的人，但他只在特殊情况下寻求整数解，通常满足于找到有理数中的一个解（II）。那是他通常能够通过代数计算解决的一类问题，其中

* 为支持这一论点，还可以指出，完全数定理的证明基本上只是唯一素因子分解定理的另一个特例。此外，所有证据表明，从这一时期起，将一个明确的数分解为素数已是众所周知并经常使用的；但在高斯于《算术研究》开头给出的证明之前，找不到分解定理的完整证明（（VIII）第I卷，第15页）。

** 参见 B. L. van der WAERDEN，《毕达哥拉斯学派的算术》，Math. Ann.，第CXX卷（1947-49），第127页。前一版本遗留段落的一个例子是第九卷第21至34命题，它们讨论被2整除的最基本性质，无疑可追溯到素数一般理论尚未发展的时期。此外，众所周知，偶数和奇数范畴在早期毕达哥拉斯学派形而上学的思辨中起了重要作用，因此自然倾向于将这一部分归于他们。

未知数的算术性质并不相关*；此外，整除性理论只起次要作用（“素数”一词仅出现一次（（II），第五卷，问题9，第1卷，第334-335页），而互素数的概念仅在涉及如下定理时才被引用：该定理断言，两个互素数的商能成为平方数，仅当它们各自都是平方数）**。

对不定方程整数解的研究，真正开始于中世纪早期的中国和印度数学家。前者似乎是由编制历法的实际问题引导到这类考虑的（其中确定各种天文现象循环的共同周期，恰恰构成一个线性的“丢番图”问题）；无论如何，他们（确定在公元4世纪至7世纪之间）提出了求解一次同余式组的法则（参见VI，第33页，习题25）。至于印度人，其数学在公元5世纪至13世纪间繁荣发展，他们不仅知道如何有条理地处理（通过应用欧几里得算法）任意多个未知数的线性丢番图方程组***，而且他们率先攻克并解决了二次问题，其中包括“费马方程”的某些特殊情形 \(Nx^{2} + 1 = y^{2}\) （（III），第II卷，第87-307页）。

此处不宜详述次数 \textgreater1 的丢番图方程理论的历史，该理论经由费马、欧拉、拉格朗日和高斯的工作，在 19 世纪导向了代数整数理论。正如我们已指出的（参见第II章和第III章的历史注记），线性方程组的研究在这一时期相当被忽视，因为它似乎不再呈现值得关注的问题：特别是，既没有寻求任意方程组的一般存在条件，也没有描述解集。尽管如此，埃尔米特在他于 19 世纪

* 虽然丢番图关于不定问题的工作总是归结为单一未知数的问题，即通过数值选取其他未知数，以使其最终方程有解，但使用这一方法的主要原因似乎是他的记号不允许他同时用多个未知数进行计算；无论如何，他在整个计算过程中都记录他所做的数值代入，并在需要时通过写出代入变量的相容性条件并先解这个辅助问题来修改它们。换言之，他处理这些代入的数值就像我们处理参数一样，以至于他实际所做的相当于求一个给定代数簇或其子簇的有理参数表示（参见（II bis））。

世纪的数论研究中，不得不使用关于线性丢番图方程的若干引理，特别是整数系数线性变换的一个“简化形式”（（XIII），第 164 页和第 265 页）；最后，在 1858 年赫格尔给出了秩等于方程个数的方程组的存在条件之后，H. J. 史密斯在 1861 年定义了整数矩阵的不变因子，并得到了如下一般定理：这样的矩阵可化为我们在 VII，第 22 页，推论 1 中给出的“标准形”（XVII）。

但与此同时，随着高斯引入这一概念（参见第一章、第二章和第三章的历史注释），以及它在数论后续发展中所起的重要作用，阿贝尔群的概念逐渐被精确化。在他于《算术研究》中对给定判别式的二次型类所构成的有限阿贝尔群所作的特别深入的研究中，高斯很快意识到其中某些群并非循环群：“在这种情况下，”他说，“一个基” {[}也就是说，一个生成元{]} 不足以满足要求，有必要取两个或更多个，通过乘法和合成*，可以产生所有类»（（VIII），第 I 卷，第 374-375 页）。不能确定高斯是否打算用这些话来描述群作为循环群的直积的分解；然而，在《算术研究》的同一篇文章中，他证明了群中存在一个元素，其阶是所有元素阶的最小公倍数——换言之，他得到了群的最大不变因子的存在性（（VIII），第 I 卷，第 373 页）；另一方面，直积的概念他是知道的，因为在一份写于1801年但生前未发表的手稿中，他勾勒了把有限阿贝尔群分解为 \(p\) -群的直积的一般证明 ** （（VIII），第 II 卷，第 226 页）。无论如何，在1868年，高斯著作集的编辑谢林，受到这些结果（特别是他刚刚发现的这份手稿）的启发，证明了（仍然针对二次型类群）一般的分解定理 (XVIII)，其方法——正如两年后克罗内克 (XX) 以抽象术语所重复的那样——本质上就是我们上面使用过的方法 (VII，第18页，定理1)。至于无挠阿贝尔群，我们已经说过（参见第II章和第III章的历史注记），由高斯、阿贝尔和雅可比发展的椭圆函数和阿贝尔积分理论，是如何逐渐引导人们关注其结构的；第一个也是最著名的无限群分解为循环群直和的例子，是1846年狄利克雷在他关于代数数域单位的论文中给出的 (XI)。但直到1879年，有限生成阿贝尔群理论与史密斯定理之间的联系才被弗罗贝尼乌斯和施蒂克尔贝格 ((XXIII)，第10节) 认识到并明确使用。

大约在同一时期，矩阵（实或复元素）的相似理论也趋于完成。线性变换特征值的概念明确出现在常系数线性微分方程组理论中，拉格朗日 (VIa) 将其应用于小振动理论，拉格朗日 (VIb) 和拉普拉斯 (VIIa) 将其应用于行星的“长期”摄动。它隐含在许多其他问题中，这些问题也在18世纪中叶左右被研究，例如寻找圆锥曲线或二次曲面的轴的问题（首先由欧拉 (Va) 解决），或者（也由欧拉 (Vb) 发展的）刚体主惯性轴的研究（由德·塞格纳于1755年发现）；我们现在知道，它（以更隐蔽的形式）也涉及偏微分方程理论的初期，特别是振动弦方程。但是（撇开最后这种情况不谈），在柯西 (X) 之前，这些问题之间的联系几乎没有被认识到。此外，由于其中大多数问题涉及对称矩阵的使用，特征值最初之所以被研究，主要是因为这些矩阵；我们将在本专著专门讨论埃尔米特算子的各章之后的历史注记中更详细地回到这一点；这里只需指出，早在1826年，柯西就证明了此类矩阵的特征值在相似变换下的不变性，并证明了它们对于 \(3 \times 3\) 对称矩阵 (Xa) 是实数，三年后 (Xb) 他将此结果推广到任意实对称矩阵。* 由莫比乌斯于1827年引入的投影的一般概念，迅速导致了对此类变换进行分类的问题（首先在2维和3维），这无非就是相应的矩阵在相似意义下的分类问题；但很长一段时间里，这个问题仅用19世纪中叶流行的“综合”方法处理，其（无论如何相当缓慢的）进展似乎并未以任何方式影响特征值理论。另一个几何问题，即圆锥曲线或二次曲面“束”的分类，情况则不同，从现代观点来看，这相当于研究矩阵 \(U + \lambda V\) ，其中 \(U\) 与 \(V\) 是两个对称矩阵；西尔维斯特在 1851 年着手解决这个问题时，肯定就是本着这种精神，仔细考察（为所考虑的束找出“标准形”）矩阵 \(U + \lambda V\) 的子式当为 \(\lambda\) 代入一个使行列式为零的值 (XIV) 时会发生什么。特征值理论的纯代数方面同时也在发展；因此，大约在 1850 年，几位作者（包括西尔维斯特本人）证明了 \(U^n\) 的特征值是 \(n\) 次幂，即 \(U\) 的特征值的幂，而凯莱在 1858 年，在他引入矩阵算术的那篇论文（XVI）中，对任意方阵*宣布了“凯莱-哈密顿定理”，尽管他只对 \(2 \times 2\) 与 \(3 \times 3\) 矩阵给出了直接证明。最终，魏尔斯特拉斯在 1868 年，利用西尔维斯特的方法，得到了“束” \(U + \lambda V\) 的“标准形”，此时 \(U\) 与 \(V\) 是方阵，不必对称，只需满足条件 \(\det(U + \lambda V)\) 不恒为零；他由此推导出任意（复）方阵的初等因子的定义，并证明了它们刻画了矩阵在相似意义下的等价类（XIX）；顺便说一句，这些结果后来由若尔当在两年后部分地（且显然是独立地）重新得到**（XXI）。这里又是弗罗贝尼乌斯在1879年指出，魏尔斯特拉斯定理可以容易地从史密斯定理（推广到多项式的情形）推出（（XXII），§ 13）；他的方法正是我们上文给出的该定理证明的基础（VII，第35页）。

我们刚才提到了单变量多项式的整除性理论；多项式的除法问题自然必定在代数学的最早期就已出现，因为它是乘法的逆运算（后者甚至丢番图就已经知道，至少对于低次多项式是如此）；但可以想象，在变量的各次幂的连贯记法建立之前，几乎不可能以一般方式处理这个问题。事实上，在我们所知的16世纪中叶之前，几乎找不到多少“欧几里得”式多项式除法程序的例子***；而S.斯蒂文（基本上使用指数记法）似乎是第一个想到将“欧几里得算法”推广以求出两个多项式的最大公因式的人（（IV，第I卷，第54-56页））。除此之外，直到18世纪中叶，整除性理论一直局限于有理整数。正是欧拉在1770年开创了算术的新篇章，他多少有些轻率地将整除性概念推广到二次扩域中的整数：在寻求确定形如 \(x^{2} + cy^{2}\) ( \(x, y\) 与 \(c\) 为有理整数）的数的因子时，他令

\[
x + y \sqrt {- c} = (p + q \sqrt {- c}) (r + s \sqrt {- c}) (p, q, r, s \text { rational   integers })
\]

并且，对两边取范数，他毫不犹豫地断言， \(x^{2} + cy^{2}\) 的所有因子都可以这样得到，形如 \(p^{2} + cq^{2}\) （Vc）。换言之，

欧拉仿佛认为环 \(Z[\sqrt{-c}]\) 是主理想整环；稍后他又用类似的论证将“无穷递降”方法应用于方程 \(x^{3} + y^{3} = z^{3}\) （他将问题归结为求 \(p^{2} + 3q^{2}\) 的立方根，为此他令 \(p + q\sqrt{-3} = (r + s\sqrt{-3})^{3}\) ）。但拉格朗日在1773年就证明了（VIc），形如 \(x^{2} + cy^{2}\) 的数的因子并非都具有这种形式，这是后来在高斯及其后继者对分圆域中整除性的研究中将更加清楚地显现出来的根本困难的第一个例子*；一般来说，不可能将有理整数整除性的基本性质（如最大公因式的存在性和素因子分解的唯一性）直接推广到这些域。这里不适合描述库默尔如何针对分圆域（XII）**，以及随后戴德金和克罗内克如何针对任意代数数域，通过理想理论的发明成功地克服了这一巨大障碍——这是现代代数中最具决定性的进展之一。但戴德金始终对各种数学理论的基础充满好奇，他并不满足于这一成功；通过分析整除性关系的机制，他在一篇论文中奠定了格序群理论的基础（这篇论文当时不为同时代人所知，湮没无闻达30年之久），它无疑是公理化代数学最早的著作之一（XXIV）；除去记号上的差异，他的工作已非常接近我们在第六章§ 1中呈现的这一理论的现代形式。

从18世纪中叶起，当时的主要任务是寻找“代数基本定理”（参见第四、五章的历史注记）。这里我们无需提及达朗贝尔的尝试，他开创了使用无穷小演算的证明系列（参见《一般拓扑学》第八章的历史注记）。但欧拉在1749年以完全不同的方式处理了这个问题

* 高斯似乎曾一度希望 \(n\) 次单位根域中的整数环会是主理想整环；在他生前未发表的一份手稿中（（VIII），第II卷，第387-397页），他证明了三次单位根域中存在欧几里得除法过程，并给出五次单位根域中类似过程的一些提示；他利用这些结果，通过比欧拉的论证更正确的“无穷递降”方法证明了方程 \(x^3 + y^3 = z^3\) 在三次单位根域中无解，并指出可以将该方法推广到方程 \(x^5 + y^5 = z^5\) ，但在方程 \(x^7 + y^7 = z^7\) 处止步，说在这里无法先验地排除 \(x, y\) 与 \(z\) 不被 7 整除的情形。

(Vd)：对每个实系数多项式 f，他试图证明存在分解 \(f = f_{1}f_{2}\) ，分解为两个（非常数的）实系数多项式，这将使他通过对 f 的次数作归纳来证明“基本定理”。正如他所注意到的，甚至只需止步于第一个奇数次因子即可，因此所有困难都归结为考虑 f 的次数 n 为偶数的情况。欧拉随后把自己限制在所求因子次数各为 n/2 的情形，并指出通过适当的消元过程，可以把 \(f_{1}\) 与 \(f_{2}\) 的未知系数表示为一个实系数方程的根的有理函数，该方程的首末两项符号相反，因而至少有一个实根。但欧拉的证明只是一个草图，跳过了许多要点；直到 1772 年，拉格朗日才通过极其冗长而细致的分析成功解决了该证明提出的困难（VId），在此过程中他展示了在运用他新近创造的“伽罗瓦方法”方面的非凡技巧（参见第四章和第五章的历史注记）。

尽管如此，拉格朗日与欧拉及其同时代的所有人一样，毫不犹豫地在多项式的“根域”内进行形式推理（也就是说，用他的语言，考虑该多项式的“虚根”）；他那个时代的数学并未为这类论证提供任何合理性证明。高斯从一开始就坚决反对18世纪无限制的形式主义，在他的学位论文中强烈抨击了这种滥用（（VIII），第III卷，第3页）。但他不可能没有察觉到，这本质上是一个内在正确论证的表面上有缺陷的表述。几年后，我们还发现他（（VIII），第III卷，第33页；另见（VIIIBis））采用了欧拉论证的一个更简单的形式，该形式由丰瑟内于1759年提出（然而，丰瑟内未能利用它取得任何成果），以获得“基本定理”的一个新证明，在该证明中他小心地避免提及所有“虚”根：后者被巧妙地添加和特殊化不定元所取代。我们在正文中给出的本质上正是高斯的这个证明（VI，第26页，定理3），并利用代数扩张进行了简化。

因此，拓扑在“基本定理”中的作用被归结为这样一个定理：具有实系数的多项式在区间内不能在不取零值的情况下改变符号（博尔扎诺多项式定理）。该定理也是所有关于（具有实系数的）多项式实根分离判据的基础，这是19世纪最受青睐的代数课题之一*。在这项研究过程中，显而易见的是，起关键作用的是R的序结构，而非其拓扑结构*；例如，博尔扎诺多项式定理在全体实代数数域中仍然成立。这一思路最终导致了由E.阿廷和O.施赖埃尔（XXV）创建的序域的抽象理论；其最显著的结果之一无疑是发现了域上序关系的存在与该域的纯代数性质相关。这就是第六章第2节所呈现的理论。

